\Article{13}
The principal task now consists in a transformation of the function \(T\) similar to the one we presented in Theoria Motus Corporum Coelestium, Art. 182, and more fully in Disquisitione de elementis ellipticis Palladis. Namely, we shall set~(17)
\[
\begin{aligned}\relax
[bb,1]&=[bb]-\frac{[ab]^2}{[aa]}\\[4pt]
[bc,1]&=[bc]-\frac{[ab][ac]}{[aa]}\\[4pt]
[bd,1]&=[bd]-\frac{[ab][ad]}{[aa]}\\
\text{etc.}&\\
[cc,2]&=[cc]-\frac{[ac]^2}{[aa]}-\frac{[bc,1]^2}{[bb,1]}\\[4pt]
[cd,2]&=[cd]-\frac{[ac][ad]}{[aa]}-\frac{[bc,1][bd,1]}{[bb,1]}\\
\text{etc.}&\\
[dd,3]&=[dd]-\frac{[ad]^2}{[aa]}-\frac{[bd,1]^2}{[bb,1]}-\frac{[cd,2]^2}{[cc,2]}
\end{aligned}
\]
\SourceLeaf{075}
etc. etc. Then, on setting\footnote{In the preceding discussion, sets of three letters could suffice for the various systems of quantities referring to the first three conditional equations; here, however, it seemed appropriate to add a fourth, so that the rule of the algorithm might appear more clearly. And since the letters \(a,b,c;\ A,B,C;\ \mathfrak A,\mathfrak B,\mathfrak C\) in their natural sequence are naturally followed by \(d,D,\mathfrak D\), while in the sequence \(x,y,z\) the alphabet is exhausted, we have appended \(w\), and likewise to the sequence \(\xi,\eta,\zeta\) we have appended \(\varphi\).}
\[
\begin{aligned}\relax
[bb,1]y+[bc,1]z+[bd,1]w+\text{etc.}&=\eta'\\
[cc,2]z.+[cd,2]w+\text{etc.}&=\zeta''\\
[dd,3]w+\text{etc.}&=\varphi'''
\end{aligned}
\]
\begin{center}etc., we shall have\end{center}
\[
T=\frac{\xi\xi}{[aa]}+\frac{\eta'\eta'}{[bb,1]}+\frac{\zeta''\zeta''}{[cc,2]}+\frac{\varphi'''\varphi'''}{[dd,3]}+\text{etc.}
\]
and the quantities \(\eta',\zeta'',\varphi'''\) etc. depend on \(\xi,\eta,\zeta,\varphi\) etc. through the following equations:
\[
\begin{aligned}
\eta'&=\eta-\frac{[ab]}{[aa]}\xi\\[5pt]
\zeta''&=\zeta-\frac{[ac]}{[aa]}\xi-\frac{[bc,1]}{[bb,1]}\eta'\\[5pt]
\varphi'''&=\varphi-\frac{[ad]}{[aa]}\xi-\frac{[bd,1]}{[bb,1]}\eta'-\frac{[cd,2]}{[cc,2]}\zeta''\\
\text{etc.}&
\end{aligned}
\]

All the formulas necessary for our purpose can now easily be obtained from this. Namely, to determine the correlates \(A,B,C\) etc., we shall set (18)
\[
\begin{aligned}
\mathfrak B'&=\mathfrak B-\frac{[ab]}{[aa]}\mathfrak A\\[5pt]
\mathfrak C''&=\mathfrak C-\frac{[ac]}{[aa]}\mathfrak A-\frac{[bc,1]}{[bb,1]}\mathfrak B'\\[5pt]
\mathfrak D'''&=\mathfrak D-\frac{[ad]}{[aa]}\mathfrak A-\frac{[bd,1]}{[bb,1]}\mathfrak B'-\frac{[cd,2]}{[cc,2]}\mathfrak C''
\end{aligned}
\]
etc.; then \(A,B,C,D\) etc. will be obtained by the following formulas, in reverse order, beginning with the last,
\[
\left.\begin{aligned}\relax
[aa]A+[ab]B+[ac]C+[ad]D+\text{etc.}&=\mathfrak A\\
[bb,1]B+[bc,1]C+[bd,1]D+\text{etc.}&=\mathfrak B'\\
[cc,2]C+[cd,2]D+\text{etc.}&=\mathfrak C''\\
[dd,3]D+\text{etc.}&=\mathfrak D'''
\end{aligned}\right\}\ (19)
\]
\begin{center}etc.\end{center}
\SourceLeaf{076}
For the sum \(S\), however, we have the new formula (20)
\[
S=\frac{\mathfrak A\mathfrak A}{[aa]}+\frac{\mathfrak B'\mathfrak B'}{[bb,1]}+\frac{\mathfrak C''\mathfrak C''}{[cc,2]}+\frac{\mathfrak D'''\mathfrak D'''}{[dd,3]}\quad\text{etc.}
\]
\Needspace{12\baselineskip}
Finally, if the weight \(P\) to be assigned to the most plausible determination of the quantity expressed by the function \(u\) is desired, we shall set (21)
\[
\begin{aligned}\relax
[bl,1]&=[bl]-\frac{[ab][al]}{[aa]}\\[5pt]
[cl,2]&=[cl]-\frac{[ac][al]}{[aa]}-\frac{[bc,1][bl,1]}{[bb,1]}\\[5pt]
[dl,3]&=[dl]-\frac{[ad][al]}{[aa]}-\frac{[bd,1][bl,1]}{[bb,1]}-\frac{[cd,2][cl,2]}{[cc,2]}
\end{aligned}
\]
etc.; this gives (22)
\[
\frac1P=[ll]-\frac{[al]^2}{[aa]}-\frac{[bl,1]^2}{[bb,1]}-\frac{[cl,2]^2}{[cc,2]}-\frac{[dl,3]^2}{[dd,3]}-\text{etc.}
\]

Formulas (17)\(.\ .\ .\ .\)(22), whose simplicity seems to leave nothing to be desired, provide a solution to our problem that is complete in every respect.

\Article{14}
After completing the principal problems, we shall also touch on some secondary questions that will shed more light on this subject.

First we must investigate whether the elimination. by which \(x,y,z\) etc. must be derived from \(\xi,\eta,\zeta\) etc. can ever become impossible. This will clearly happen if the functions \(\xi,\eta,\zeta\) etc. are not independent of one another. Let us therefore suppose for the moment that one of them is already determined by the others, so that we have the identity
\[
\alpha\xi+\beta\eta+\gamma\zeta+\text{etc.}=0
\]
where \(\alpha,\beta,\gamma\) etc. denote determinate numbers. We shall therefore have
\[
\begin{aligned}
\alpha[aa]+\beta[ab]+\gamma[ac]+\text{etc.}&=0\\
\alpha[ab]+\beta[bb]+\gamma[bc]+\text{etc.}&=0\\
\alpha[ac]+\beta[bc]+\gamma[cc]+\text{etc.}&=0
\end{aligned}
\]
\Needspace{23\baselineskip}
etc.; hence, if we set
\SourceLeaf{077}
\[
\begin{aligned}
\alpha a+\beta b+\gamma c+\text{etc.}&=p\theta\\
\alpha a'+\beta b'+\gamma c'+\text{etc.}&=p'\theta'\\
\alpha a''+\beta b''+\gamma c''+\text{etc.}&=p''\theta''
\end{aligned}
\]
etc., it follows directly that
\[
\begin{aligned}
a\theta+a'\theta'+a''\theta''+\text{etc.}&=0\\
b\theta+b'\theta'+b''\theta''+\text{etc.}&=0\\
c\theta+c'\theta'+c''\theta''+\text{etc.}&=0
\end{aligned}
\]
etc., and also
\[
p\theta\theta+p'\theta'\theta'+p''\theta''\theta''+\text{etc.}=0
\]
Since all \(p,p',p''\) etc. are by their nature positive quantities, this equation clearly cannot hold unless \(\theta=0,\ \theta'=0,\ \theta''=0\) etc.

Let us now consider the values of the total differentials \(\mathrm dX,\mathrm dY,\mathrm dZ\) etc. corresponding to those values of the quantities \(v,v',v''\) etc. to which the observations relate. By the conclusion we have just reached, these differentials, namely
\[
\begin{aligned}
&a\mathrm dv+a'\mathrm dv'+a''\mathrm dv''+\text{etc.}\\
&b\mathrm dv+b'\mathrm dv'+b''\mathrm dv''+\text{etc.}\\
&c\mathrm dv+c'\mathrm dv'+c''\mathrm dv''+\text{etc.}
\end{aligned}
\]
etc., will depend on one another in such a way that, when multiplied respectively by \(\alpha,\beta,\gamma\) etc., they produce a sum that vanishes identically; or, equivalently, any one of them whose corresponding multiplier \(\alpha,\beta,\gamma\) etc. is nonzero will itself vanish as soon as all the others are assumed to vanish. Consequently, at least one of the conditional equations \(X=0,\ Y=0,\ Z=0\) etc. must be regarded as \textit{redundant}, since it is satisfied automatically as soon as the others are satisfied.

However, if we look more deeply into the matter, it appears that this conclusion by itself holds only for an infinitely small range of variation of the indeterminates. More precisely, two cases must be distinguished: one in which one of the conditional equations \(X=0,\ Y=0,\ Z=0\) etc. is already contained in the others absolutely and generally, which can easily be avoided in any case; the other in which, as if by chance, for those concrete values of the quantities \(v,v',v''\) etc. to which the observations re\SourceLeaf{078}%
late, one of the functions \(X,Y,Z\) etc., for example the first, \(X\), acquires a maximum or minimum value (or, more generally, a stationary value) with respect to all the changes we could make to the quantities \(v,v',v''\) etc. while keeping the equations \(Y=0,\ Z=0\) etc. satisfied. Yet since in our investigation the variation of the quantities is considered only within limits so narrow that it can be treated as infinitely small, this second case (which will hardly ever occur in practice) will have the same effect as the first: namely, one of the conditional equations will have to be rejected as redundant. We can be certain that, if all the retained conditional equations are independent of one another in the sense intended here, elimination will necessarily be possible. However, we must reserve a fuller investigation for another occasion; this subject is not undeserving of one on account of its theoretical subtlety rather than its practical usefulness.

\clearpage
